\section{Restricted lattices, auxiliary stages, and coordinate pins}
\label{sec:stage-existence}

\status{SE1--SE16 prove the scalar-descent and restricted-lattice
construction of the literal divisibility stage, the finite-dimensional
row model, the ordinary local and spectral-stage model, the compact-dual
Haar normalization and the jet-stage model, and coordinate normalization
of the two pins. The precise remaining condition in the assembled
\texttt{Later} constructor is an \texttt{NStage} on the same adapted row
object. Its local period-ring and Galois suppliers are not declared
proved here. These existence statements do not assert the quantitative
spectral, metric, cohomological or prime-window properties needed later.}

In SE1--SE6, $K$ is a number field when integral lattices or infinite
places are used, $\sigma:K\hookrightarrow\mathbb C$ is fixed, and a
$\sigma$-semilinear map satisfies $e(ax)=\sigma(a)e(x)$.
For a Hermitian lattice $A=(V,b,X,L)$, our convention is
\[
 \|x\|_{A,v}=\|X_v^{-1}\sigma_v(b^{-1}x)\|_2.
 \tag{SE.1}
\]
The word ``good'' means that every $X_v$ is invertible, $L$ is a finitely
generated $\mathcal O_K$-module, and $KL=V$. It does not assert that $L$
is free over $\mathcal O_K$.

\paragraph{SE1: scalar extension preserves every independent family.}
Let $b_1,\ldots,b_n$ be a $K$-basis of $V$, and suppose
$e(b_1),\ldots,e(b_n)$ are complex-linearly independent in $M$.
Put $M_0=\operatorname{span}_{\mathbb C}e(V)$. Every semilinear map
$g:V\to N$ extends to a complex-linear map on $M$: prescribe
$e(b_i)\mapsto g(b_i)$ on $M_0$ and extend by zero on a vector-space
complement. This gives $F(e(x))=g(x)$ because
$e(x)=\sum_i\sigma(b_i^*(x))e(b_i)$.

In particular, $e$ is injective. If $z_1,\ldots,z_t$ are independent over
$K$, extend their coordinate functionals from their span to $V$ and put
$g(x)=(\sigma(z_1^*(x)),\ldots,\sigma(z_t^*(x)))$.
An extension $F$ sends $e(z_j)$ to the $j$th standard vector. Applying
$F$ to a complex linear relation proves that the $e(z_j)$ are independent
over $\mathbb C$. This argument does not require $\sigma(K)$ to be
topologically dense, or $e(V)$ to be a complex subspace.

\paragraph{SE2: a finite common kernel descends to the coefficient field.}
Let $T_s\in\operatorname{End}_{\mathbb C}(M)$ for $s$ in a finite set
$S$, and assume $T_s(e(V))\subseteq e(V)$ for every $s$. Then
\[
 M_0\cap\bigcap_{s\in S}\ker T_s
 =\operatorname{span}_{\mathbb C}
 e\left(\{x\in V:T_s e(x)=0\text{ for all }s\}\right).
 \tag{SE.2}
\]
Injectivity in SE1 gives unique $K$-linear $N_s$ with $eN_s=T_se$.
Combine these into $N:V\to V^S$, and choose a right inverse of
$N:V\to\operatorname{im}N$. Extend that right inverse to a linear
$r:V^S\to V$. The map $p=1-rN$ takes values in $\ker N$ and is the
identity there. If $\iota_s:V\to V^S$ is the $s$th coordinate inclusion,
\[
 1_V=p+\sum_s r\iota_sN_s.
 \tag{SE.3}
\]
Extend $ep$ and each $er\iota_s$ by SE1 to complex-linear maps $P,R_s$.
Equation (SE.3), first on the $e(b_i)$ and then on their complex span,
gives $1_{M_0}=P+\sum_sR_sT_s$. Also $P(M_0)$ is contained in the
right side of (SE.2). Thus a common-kernel vector $z\in M_0$ equals
$Pz$ and belongs to that right side. The reverse inclusion follows from
complex linearity of the $T_s$. The finite index set is essential to
this displayed finite-sum proof; no assertion about arbitrary infinite
intersections is being substituted for it.

\paragraph{SE3: the restricted metric has explicit frames.}
Let $U\subseteq V$ be a $K$-subspace with basis $c_1,\ldots,c_t$. For
each infinite place $v$, let $E_v$ be the matrix whose columns are
$\sigma_v(b^{-1}c_j)$. By SE1, $E_v$ has full column rank. Hence
\[
 A_v=(X_v^{-1}E_v)^*(X_v^{-1}E_v)
 \tag{SE.4}
\]
is positive definite. A positive definite Hermitian matrix is unitarily
diagonalizable with positive eigenvalues; taking their positive square
roots gives an invertible $B_v$ with $B_v^*B_v=A_v$. Set $Y_v=B_v^{-1}$.
For $x=\sum_j a_jc_j$,
\[
 \|Y_v^{-1}\sigma_v(a)\|_2^2
 =\sigma_v(a)^*A_v\sigma_v(a)
 =\|X_v^{-1}\sigma_v(b^{-1}x)\|_2^2.
 \tag{SE.5}
\]
Taking nonnegative square roots proves exact equality of norms, at
every infinite place. The empty matrix treats $U=0$: its determinant
is $1$ and both norms are zero. No numerical approximation to a matrix
square root enters the statement.

\paragraph{SE4: intersection with a rational subspace is a full lattice.}
Give $U$ the lattice
$L_U=\{x\in U:x\in L\}$ and the frames of SE3. Its image in $V$ is
exactly $L\cap U$. Since $\mathcal O_K$ is Noetherian, $L\cap U$ is a
finitely generated submodule of the finitely generated module $L$.
For $x\in U$, write $x$ as a finite $K$-linear combination of vectors
of $L$. Clear all denominators by some
$0\ne a\in\mathcal O_K$. Then $ax\in L\cap U$ and
$x=a^{-1}(ax)$ belongs to $KL_U$. Thus $L_U$ spans $U$, and together
with SE3 it is a good Hermitian lattice. This proof uses neither a
principal-ideal assumption nor a global integral basis for $L_U$.
The standard lattice $\mathcal O_K^d\subseteq K^d$ with identity
frames is good for the same reasons, including $d=0$.

\paragraph{SE5: the ramified subspace is the extension of its actual form.}
Fix the source's member forms $M$ and column data $A$. Its embedding
$e=M.\mathrm{emb}(0)$ has complex-independent basis image spanning
$A.\mathrm{Mu0}$, by \texttt{emb\_basis} and \texttt{Fi\_zero}.
The operators $\rho_v(a)$ preserve $e(V)$: they preserve the embedded
integral lattice by \texttt{rhoT\_lat}, and every vector is a
$K$-linear combination of lattice vectors. Their finite sums
$\mathrm{Nm}_v=\sum_a\rho_v(a)$ therefore also preserve $e(V)$.

Define
\[
 U=\{x\in V:e(x)\in A.\mathrm{Mram}\},\qquad
 A.\mathrm{Mram}=A.\mathrm{Mu0}\cap\bigcap_v\ker\mathrm{Nm}_v.
 \tag{SE.6}
\]
The local-prime index set is finite. SE2 proves that $e(U)$ spans this
exact complex ramified subspace. SE1 proves that every $K$-basis of $U$
maps to a complex-independent family there. Consequently the source's
\texttt{ramEmb} is a scalar-extension presentation of the actual
\texttt{Mram}; equality of dimensions alone is not used as an
identification. Its lattice is the actual intersection
$M.\mathrm{ramL}=M.\mathrm{lat}(0).L\cap U$ from SE4.

\paragraph{SE6: the literal divisibility-stage record exists.}
For every $B,C$ and every $W:\texttt{RowsRest}(B,C)$ there is a
$V:\texttt{DivStage}(B,C,W)$ whose lattices are all good and such that
\[
 \nu_u\bigl(\mathrm{Nu}(u).L\bigr)=(C.\mathrm{fm}(u)).\mathrm{ramL}.
 \tag{SE.7}
\]
Take $\mathrm{RR}=K^{D_{\rm rows}}$ with its standard lattice and
identity frames. For each $u$, take $\mathrm{Nu}(u)$ to be precisely
the subspace, intersection lattice and restricted frames of SE3--SE5.
Let $j_{EE}$ map standard coordinates through $W.\mathrm{onb}$, let
$\nu_u$ be the subspace inclusion, and set
\[
 j_{N,u}=(W.e\mathrm{Mu}_u)^{-1}\circ\mathrm{ramEmb}_u.
 \tag{SE.8}
\]
The full field check is as follows. The fields \texttt{div\_good},
\texttt{jEE\_norm} and \texttt{nuIn\_norm} follow respectively from
invertible frames, orthonormal coordinates and (SE.5).
\texttt{jEE\_basis} is the orthonormal basis assertion;
\texttt{jN\_basis} is SE5 transported through the displayed linear
equivalence. The fields \texttt{nuIn}, \texttt{nuIn\_inj} and
\texttt{nuIn\_range} are the inclusion, its injectivity and its range.
\texttt{nuIn\_jN} is (SE.8) followed by $W.e\mathrm{Mu}_u$.
Use exactly $W.D_{\rm all}$ and $W.\mathrm{delp}$ for
\texttt{Dall\_eq} and \texttt{delP\_eq}. The \texttt{div} data are
completed by
\[
 B_k=0,\quad h=0,\quad\eta_F=0,\quad
 HH=s\mathrm{Frame}(\mathrm{Nu},0)+\mathrm{RR}.\mu_{\max}.
 \tag{SE.9}
\]
Then \texttt{HH\_ge} is equality, and the bound on \texttt{delP}
is the bound already in $W$. SE4 gives \texttt{AllGood} and (SE.7).
This accounts for all fourteen \texttt{DivStage} fields. The zero
$B_k$ in this existence construction does not provide a later estimate
for independently prescribed nonzero maps.

\paragraph{SE7: the actual finite class data supply \texttt{Fits}.}
For an actual period datum $Q$ with $Q.\mathrm{Ok}$ and $Q.H(b)$ the
full group, evaluate an equivariant form on one representative of every
class. This defines a linear map
\[
 Q.M_t\longrightarrow\prod_{c\in Q.S'.\mathrm{Cl}} VV_c.
 \tag{SE.10}
\]
The class set and the polynomial representation spaces $VV_c$ are
finite. If all evaluations vanish, equivariance and the specified
movement from the representative to any point force all values to
vanish; thus (SE.10) is injective and $Q.M_t$ is finite dimensional.

Now fix the actual $I$ and $l:I.\mathrm{Laws}$, and put $C=I.\mathrm{cols}(l)$.
The source period spaces are independent of the word $u$ at fixed
twist $\mu$, with their stated common polynomial type. Set
$d(\mu)=\dim_{\mathbb C}(I.\mathrm{perIn}(u_0,\mu).M_t)$.
Choose coordinate linear equivalences to $\mathbb C^{d(\mu)}$;
these give both \texttt{mt0} and \texttt{mtTw} on this same $C$.
By TL7 (or the weighted positive class-sum calculation FC7),
\[
 \operatorname{ip}_u(x,y)=
 \sum_c w_c\langle x(c),y(c)\rangle_c
 \tag{SE.11}
\]
is conjugate-linear in $x$, linear in $y$, Hermitian, and positive
definite on $\mathrm{Mu0}$: every $w_c>0$, and a nonzero equivariant
form has a nonzero value on some representative. Restriction to
$\mathrm{Mram}$ gives \texttt{herm}, \texttt{pos}, \texttt{defn}.
All five fields of $\texttt{Fits}(C,d)$ are supplied. We retain the
literal $I.\mathrm{Laws}$ and finite-class hypotheses.

\paragraph{SE8: the row record has an explicit finite-dimensional model.}
For every $C$,
\[
 \texttt{Nonempty}\,\texttt{RowsRest}(B,C)
 \quad\Longleftrightarrow\quad
 \exists d:\ C.\mathrm{Tup}\to\mathbb N,\ \texttt{Fits}(C,d).
 \tag{SE.12}
\]
For the forward implication use $d(\mu)=\dim W.\mathrm{Amb}(\mu)$,
the two equivalences $W.eM_{t,0},W.eM_{t,\mu}$, and the isometry
$W.e\mathrm{Mu}$ to recover the five fields of \texttt{Fits}.

For the reverse implication choose arbitrary nonnegative $n,n_0,a$.
Use $EE=\mathbb C^n$, $EE_0=\mathbb C^{n_0}$ with their standard
orthonormal bases, $D_{\rm all}=2n$, $D_{\rm rows}=n$,
$\mathrm{delp}=1/2$, and a singleton $FT$. Set
$\mathrm{Amb}(\mu)=\mathbb C^{d(\mu)}$ and
$\mathrm{Mult}(\mu,0)=\mathbb C^n$, with all positive-level
$\mathrm{Mult}(\mu,j)=\mathbb C^{d(\mu)}$.
Every transition $io$ is coordinate truncation or extension by zero;
the top identification uses $m\ge2$.
Set $\mu\mathrm{At}=\mu_0$ and $\mu\mathrm{Tw}=\mu_\star$.

Take $\mathrm{Mu}(u)$ to be the actual $\mathrm{Mram}$, endowed with
(SE.11), and $e\mathrm{Mu}$ the identity. Transfer the actual base
period $P\widetilde{\mathrm{Col}}$ along the chosen \texttt{mt0}
equivalence; for other twists transfer the actual twisted period along
\texttt{mtTw}, including the given $e0\mathrm{Tw}$ identification.
Define $B_u$ by the base period after coordinate truncation
$\mathbb C^n\to\mathbb C^{d(\mu_0)}$. Its two linearity fields
hold because the finite column chain is a composite of linear maps
and $P\widetilde{\mathrm{Col}}$ is bilinear. This is a complex
bilinear period pairing, whereas (SE.11) is Hermitian.

For the remaining fields take $M_j=M_t=\mathbb C^a$,
$M_{t,K}=K^a$, the coordinate integral lattices, and
$jM_t(x)=(\sigma(x_i))_i$. SE1, together with the spanning standard
vectors, proves \texttt{jMt\_basis} for every $K$-basis; the generators
are images of the standard $K$-vectors, proving \texttt{MtLat\_sub}.
Put $e\mathrm{Op}=eM_t=\theta_P=1$, $\mathrm{isot}=\top$,
$\lambda_P=\mathrm{ze}=1$, $\kappa_{\rm big}=1$,
$\kappa_{\rm small}=-1$, $I_j=jw=\eta_J=0$, and
$\mathrm{shaLen}=1$. Idempotence, self-adjointness, coefficient-field
preservation and nonnegativity now follow directly. The base/twist
period ties hold by the definitions above. These assignments account
for every scalar, operator, lattice, period and metric field of the
record; choosing $n=n_0=a=1$ proves the required existence.

\paragraph{SE9: the ordinary local data and the specified split places.}
Let $E/F$ be the quadratic CM extension in $B$, with its nontrivial
involution $c$, and let $v$ be a finite place of $F$ of residue
characteristic $\ell$. Take
\[
 F_v=\widehat F_v,\qquad E_v=F_v\otimes_F E,
 \quad c_v=1\otimes c,
 \quad O_v=\{x\in E_v:x\text{ integral over }\mathcal O_{F_v}\}.
 \tag{SE.13}
\]
The completion is a characteristic-zero local field: an integral basis
and the $v$-adic norm make it a finite extension of $\mathbb Q_\ell$.
The original field is dense, its normalized valuation is the given
one, and tensoring a basis of $E/F$ gives the stated free rank two and
spanning property. Give $E_v$ its finite-dimensional vector topology
and Borel sigma algebra. The standard character of $\mathbb Q_\ell$
is $x\mapsto\exp(2\pi i\{x\}_\ell)$, where the principal part
$\{x\}_\ell\in\mathbb Z[1/\ell]/\mathbb Z$ is determined by the
finitely many negative digits. Composing with
$\operatorname{Tr}_{E_v/\mathbb Q_\ell}$ gives exactly the character
and principal-part specification in \texttt{LocAt}. These choices
supply all its fields before the optional split datum.

For clarity, the local-field facts just used are ordinary completion
and tensor-base-change facts: $F_v\otimes_FE\simeq\prod_{w\mid v}E_w$,
and $[E_w:\mathbb Q_\ell]=e_wf_w$. They follow by completing an
integral basis, decomposing the completed integral algebra into its
local factors by CRT, and taking fraction fields. If $\ell$ splits
completely in $E$, the equality
$\sum_{w\mid\ell}e_wf_w=[E:\mathbb Q]$, together with the maximal
number of distinct primes, forces every $e_w=f_w=1$. Hence
$F_v=\mathbb Q_\ell$ and $E_v\simeq F_v\times F_v$; $c_v$ swaps
the two factors. This is the required split algebra and involution,
and $\ell$ is its specified uniformizer. Equivalently, for
$E=F(\alpha)$ with $\alpha^2=a$ and $c\alpha=-\alpha$, the two local
roots $\pm\beta$ give $x+y\alpha\mapsto(x+y\beta,x-y\beta)$.

Use \texttt{some} of this split datum exactly when
$\ell\in\{g_0,g_1\}$, and \texttt{none} elsewhere. The host's full
splitting hypotheses supply the required implication at these two
primes. Thus the source's biconditional for the optional datum is
satisfied. It makes no claim that other places fail to split.
This ordinary completion construction is distinct from
\texttt{LocStage}, which also requires Fontaine period-ring data.

\paragraph{SE10: the literal spectral-stage record exists.}
Suppose $B$ has its given positive gap $\Delta$, injective Boolean
letter map and exact $u$- and $v$-weight identities
(\texttt{HasLetters}); suppose (SE.11) is positive definite on
each $\mathrm{Mu0}$; and let $W$ satisfy
$D_{\rm rows}\le D_{\rm all}$. Use $HH=\mathbb C^{D_{\rm all}}$,
and give $Cu(u)=\mathrm{Mu0}(u)$ precisely the metric (SE.11).
If $b_i=W.\mathrm{onb}_i$, the map
\[
 x\longmapsto\sum_{i<D_{\rm rows}}\langle x,b_i\rangle e_i
 \tag{SE.14}
\]
is conjugate-linear under the source convention and norm preserving:
its squared norm is $\sum_i|\langle x,b_i\rangle|^2=\|x\|^2$.
It supplies the field \texttt{sEE}, including the correct semilinearity.

Choose the local data of SE9, $\mathrm{dsig}=\mathrm{cdScale}
=\mathrm{volN}=\mathrm{volH}=1$, singleton \texttt{Term} and
\texttt{Pair}, and the one-block point partition $[(1,r_S+1)]$.
Take $P_u,a\mathrm{Term},G\mathrm{term},s_0,p\mathrm{val},
\mathrm{meanCount}$ to be zero, $c\mathrm{Term}=\mathrm{massPair}=1$,
the predicates \texttt{isId}, \texttt{inS}, \texttt{relevant},
\texttt{FeqF}, \texttt{scalarPair} to be true, and
\texttt{BBne}, \texttt{generic} to be false. Use the actual
$\Delta$, letter map and weight identities; they are not replaced
by zero data. Then \texttt{s\_split} is SE9, \texttt{s\_HH} is its
dimension, \texttt{s\_Cu} and \texttt{s\_Cu\_isom} are the identity
with the chosen metric, \texttt{s\_PlF} is the identity on finite
places, and \texttt{dsig\_pos} is $1>0$. Together with (SE.14) and
the displayed \texttt{s} this checks all nine \texttt{SStage} fields.
In particular both term types are nonempty. This construction imposes
no nonzero spectral contribution.

\paragraph{SE11: compact-dual Haar normalization is constructible.}
Put $J=\operatorname{diag}(1_P,-1_Q)$ and
$\mathfrak u_{P,Q}=\{X:X^*J+JX=0\}$. The real-linear map
\[
 \mathfrak u_{P,Q}\times\mathfrak u_{P,Q}\longrightarrow M_{P+Q}(\mathbb C),
 \qquad(X,Y)\longmapsto X+iY
 \tag{SE.15}
\]
is an isomorphism. Its inverse sends $Z$ to
\[
 X=\tfrac12(Z-JZ^*J),\qquad
 Y=-\tfrac i2(Z+JZ^*J).
 \tag{SE.16}
\]
Both satisfy the defining Lie-algebra equation, and $X+iY=Z$.
Start with any Lebesgue measure $\lambda_0$ on this real vector space.
Its product pushed through (SE.15) is $a$ times entrywise complex
Lebesgue measure for some finite $a>0$. Thus
$\lambda=a^{-1/2}\lambda_0$ satisfies the literal
\texttt{IsOmegaLeb} product-measure normalization.

Here is the density fact needed to normalize the group measures. The
matrix exponential maps a neighborhood of $0$ in
$\mathfrak u_{P,Q}$ analytically onto a neighborhood of $1$ in
$U(P,Q)$ and has derivative the identity. One can see this without a
global structural theorem: the convergent local matrix logarithm is
inverse to the exponential, and $g^*Jg=J$ implies
$(\log g)^*=-J(\log g)J$. Translate a positive volume density at the
identity by left multiplication. It defines a smooth left-invariant
measure, finite on compact sets and positive on open sets. Haar
uniqueness identifies every Haar measure with a positive constant
multiple of it. In exponential coordinates its density relative to
$\lambda$ is continuous and has a finite positive value $c$ at $0$.
For the exact entrywise-norm balls
$B_r=\{X:\sum_{ij}|X_{ij}|<r\}$, these coordinates give
\[
 \frac{\mu(\exp B_r)}{\lambda(B_r)}\longrightarrow c\in(0,\infty).
 \tag{SE.17}
\]
Indeed, for any $\epsilon>0$ the continuous density lies between
$c-\epsilon$ and $c+\epsilon$ on all sufficiently small $B_r$;
integration and division by their finite positive volumes prove the
limit. The matrix logarithm also gives injectivity on those balls.
For $P+Q=0$, use the unit point masses; the ratio is identically $1$.

The compact group $U(N)$ is a closed bounded matrix set. Normalize its
Haar measure $p_N$ to total mass $1$, and let $c_N>0$ be its density
in (SE.17). For $Q>0$, multiply any Haar measure on $U(P,Q)$, of
density $d_{P,Q}>0$, by $c_{P+Q}/d_{P,Q}$. For $Q=0$ keep exactly
$p_P$. This yields all four \texttt{CDHaar} fields: a measure at
each signature, Haar invariance, probability in the compact case,
and the same finite nonzero Omega density for $(P,Q)$ and $(P+Q,0)$.
The ordinary finite-dimensional change-of-variables and Haar
existence/uniqueness theorems used in this proof are explicit classical
background; the common-density conclusion is proved here.

\paragraph{SE12: a jet-stage model with all prescribed ranks.}
Given a positive volume function, a \texttt{CDHaar}, a
\texttt{ThetaJet}, the positive \texttt{GlobalExtra} scalars, and
nonempty adjacent class sets for each of its situations, the literal
\texttt{GJStage} exists. For each situation choose the word and its
two classes. A prime of $\mathcal O_K$ above the situation's rational
prime exists by lying over, and its residue characteristic is that
prime. A fixed element of the nonempty $\mathrm{Tau}$ supplies each
required archimedean label.

For every required $u,w$, let
$d=\dim_{\mathbb C}(C.\mathrm{mem}(u)).A.\mathrm{Mu0}$ and
$\ell=\mathrm{resChar}(w)$. Use $\mathbb Z_\ell\subset\mathbb Q_\ell$
as the coefficient DVR, no variable places, the one-element level
group acting trivially, the trivial character, and $d$ components.
A polynomial in no variables is a constant. Thus the integral
automorphic module in each component is $\mathbb Z_\ell$ and the
whole module is $\mathbb Z_\ell^d$; its constant-coefficient map is
an isomorphism. The class set is a singleton in each component, so
\texttt{FinClasses} holds, and extension to $\mathbb Q_\ell$ gives
exact rank $d$. This proves the source's prescribed-rank
\texttt{GAtP} supplier, including rank zero. These constructions
fill \texttt{gat}, \texttt{gat\_d}, all jet-place and class fields;
the other fields are the supplied volumes, Haar and extra scalars.

In particular choose volume $1$, all extra scalars $1$, both situation
types empty, and the \texttt{CDHaar} of SE11. No adjacent classes or
inert-prime choices are then required. This proves
$\texttt{Nonempty}\,\texttt{GJStage}(B,C)$ for every $B,C$.
The optional source conclusion using its specific
$\mathrm{baseVol}=\mathrm{trueVol}$ additionally needs the positivity
and identification of that specific volume function; it is not
inferred from the constant-volume model. Likewise nonempty jets and
actual lift estimates require their stated additional data.

\paragraph{SE13: zero idempotents and the exact adapted row object.}
Let $r\ge2$. At each selected prime, work over its finite residue field
$\mathbb F_{g_k}$. Among monic degree-$r$ polynomials over that field
with nonzero constant term, distinguish $f_0=(X-1)^r$.
On any complex vector space define
\[
 e_f=\begin{cases}1&f=f_0,\\0&f\ne f_0.\end{cases}
 \tag{SE.18}
\]
These are central orthogonal idempotents, their support is finite,
and their sum is $1$. A good polynomial in the source's good tuple
is squarefree; $f_0$ is not squarefree since $r\ge2$. Hence every
good-polynomial idempotent is zero. The set of chosen split primes
is nonempty ($f>0$), so the product for a good tuple is zero, and
the sum $e_0$ over good tuples is zero. This proves the exact
idempotent calculation used by
\texttt{idempotent\_good\_tuples\_eq\_zero}.

The source's \texttt{adapt} is a change of row object:
\[
 e\mathrm{Op}=0,\quad eM_t=0,\quad\eta_J=0,
 \quad\mathrm{shaLen}(p,e)=
 \begin{cases}(\operatorname{shaLen}(L_p,e).\mathrm{toNat}:\mathbb Q)
                       &p\text{ prime},\\0&\text{otherwise}.
 \end{cases}
 \tag{SE.19}
\]
It preserves the other fields. Zero is self-adjoint, idempotent and
preserves the coefficient-field image; the new Sha scalar is
nonnegative by its definition. Thus these updates preserve the
\texttt{RowsRest} type. When the source's trivial central actions are
used, (SE.18) makes the \texttt{IdempotentMatchesAction} condition
hold for these zero operators. No finiteness of the Sha group follows
from (SE.19): extended-natural infinity is also sent to $0$ by
\texttt{toNat}. This lemma proves the idempotent and adaptation steps;
it does not construct the full \texttt{NStage} or its local period rings.

\paragraph{SE14: the five-stage join retains one precise input.}
Fix $I$, $l:I.\mathrm{Laws}$ and \texttt{HasLetters} for the same
$B$. Let
\[
 C=I.\mathrm{cols}(l),\qquad
 W_0=\mathrm{RowsRestModel.model}(\mathrm{fits}(l),1,1,1),
 \qquad W=\mathrm{adapt}(W_0).
 \tag{SE.20}
\]
Here the Sha scalar used in \texttt{adapt} is supplied as in SE13;
its defining cyclotomic/Galois data remain explicit if needed.
If $N:\texttt{NStage}(B,C,W)$ is supplied, then
$\texttt{Later}(B,C)$ is inhabited. Indeed, SE7--SE8 supply $W_0$,
SE13 supplies $W$ with $D_{\rm rows}=1,D_{\rm all}=2$,
SE6 supplies $V:\texttt{DivStage}(B,C,W)$, and SE10 supplies
$S:\texttt{SStage}(B,C,W)$ using $1\le2$ and the same class metric.
SE11--SE12 supply $G:\texttt{GJStage}(B,C)$. The tuple
$(W,V,S,N,G)$ is exactly the five-field \texttt{Later} record.

The dependency on this same $W$ matters. An $N$ for a different row
object cannot be substituted by a dimension comparison. The source's
remaining \texttt{exists\_nStage\_adapt} supplies such an $N$ from
\texttt{LocStage} and its other constructions; only the idempotent
part of that supplier has been proved above. In particular the
\texttt{LocStage} fields include actual $A_{\inf}$, $B_{\rm cris}$,
$B_{\rm dR}$, actions, filtration, and tied comparison data, far
beyond SE9's ordinary completions. These remain named targets.

There is a concrete reason to keep the later quantitative hypotheses
separate: in (SE.20), $M_t=\mathbb C$ but $eM_t=0$, so
$\operatorname{rank}(eM_t)=0$ on a one-dimensional space. Thus the
bare row record and its adaptation cannot imply any positive lower
bound for that projector rank. A later theorem with such a conclusion
must use further premises. The same caution applies to the zero
spectral terms in SE10 and the empty jets in SE12.

There is also a concrete arithmetic compatibility obstruction for this
particular bare-stage model. The source's \texttt{nTw} sets
$\mathrm{Tw}=\varnothing$. But the \texttt{NChoice} record requires
\[
       \mathrm{twI}:B.\mathrm{places}.\mathrm{Pl}\longrightarrow
                         \mathrm{Tw}
 \tag{SE.24}
\]
at every coefficient place, including places where $\mu\mathrm{At}=\mu_0$.
The domain is nonempty: lying over supplies a prime of the coefficient
number field above $2$, as in SE12. Thus (SE.24) cannot exist for this
specific empty-twist model. Equivalently its \texttt{ReadElsewhere}
condition fails even though its nonreference support is empty. This
does not rule out other compatible \texttt{NStage}/\texttt{Later}
constructions; it proves that the displayed bare construction cannot
by itself supply NC's arithmetic choice interface.

\paragraph{SE15: diagonal coordinates supply the auxiliary frames.}
Suppose the source ranks are fixed, $E$ has its CM involution, $T$ is
prime and $p_z\in E$ satisfies $\overline{p_z}p_z=T$. Take
$X_j=E^{N_j}$ and $X'=E^{\mathrm{rank}X'}$ with positive standard
forms and coordinate inclusions. For $j<m$ put
$h=\lfloor(m-1-j)/2\rfloor$. The rank and positive index of $Y_j$
are $N_j+(m-1-j)$ and $N_j+h$. Its first $N_j$ coordinates are
$X_j$; the next $h$ positive and $h$ negative coordinates give
\[
 u_i=e_i+f_i,\qquad v_i=(e_i-f_i)/2,
 \quad\langle u_i,u_k\rangle=\langle v_i,v_k\rangle=0,
 \quad\langle u_i,v_k\rangle=\delta_{ik}.
 \tag{SE.21}
\]
They are orthogonal to $X_j$. If an extra negative coordinate is
present, put $p_z$ there: its norm is $-T$ and it is orthogonal to all
these vectors. At the top, put $z=p_z e_{\mathrm{rank}Y_m}$ in the
extra positive coordinate; it has norm $T$ and is orthogonal to the
embedded $Y_m$. All vectors and orthogonality identities are obtained
by multiplying diagonal entries. This proves the actual
\texttt{spacesNorm\_hasFrame}, coordinate-form identities, top norm,
signatures, $V\mathrm{neg}=\mathrm{false}$, and $T$ preservation.
For completeness the remaining coordinate spaces are
$W_j=E^r$ with diagonal signature $(k+j-1,m-j)$ for $1\le j\le m$,
and $V_j=E^{n_V}$ with the positive diagonal form. Their bases and
inclusions are the standard coordinate ones, completing the
\texttt{spacesNorm\_coord} assertions.

For any $n,N$ and matching coordinate indices $a,b$, extension by
zero satisfies $\mathrm{extTo}_{n,N}(e_a)=e_b$: evaluate each
coordinate, where both sides are $1$ at the matching index and $0$
elsewhere. Applying this to $X'\hookrightarrow X_m$ is precisely
the \texttt{Pins.iX} equation, with its fixed frame indices.

\paragraph{SE16: normalizing an existing input supplies both pins.}
Start with the source's already supplied $P$, $X:\texttt{PreX}(P)$,
oriented frame data, letter identities and the required size/
$du_\star$ properties. Write $\pi_u=X.\pi_T(u)$; then
$\pi_u\overline{\pi_u}=T_u$ and every $\pi_u\ne0$, since $T_u$
is prime. Apply SE15 with $T=T_{u_0}$ and $p_z=\pi_{u_0}$, and
replace the space fields by these diagonal spaces. Define
\[
 c(u)=\pi_{u_0}/\pi_u.
 \tag{SE.22}
\]
It is nonzero, satisfies
$c(u)\overline{c(u)}=T_{u_0}/T_u$, and, by direct cancellation,
\[
 c(u)\pi_u=\pi_{u_0}.
 \tag{SE.23}
\]
Thus (SE.23) is \texttt{Pins.cof}; the coordinate computation of
SE15 is \texttt{Pins.iX}. The preserved number-field, prime, weight
and word fields transport $X$, its orientation, the letter identities,
and the stated $du_\star$ equalities exactly as in
\texttt{ExistsOkPins.exists\_pre}; all changed space fields have
been checked in SE15. The scalar norm relation above is the changed
\texttt{cOf\_eq} field. Subsequent coefficient-field and base
extension can carry these coordinate equalities unchanged.

This is a normal-form construction for an existing arithmetic and
geometric input. The existence of the initial sized
\texttt{PreX}/oriented input, its analytic prime windows, and the
later choice of compatible $I.\mathrm{Ok}$ are not consequences of
the two cancellation and coordinate equations. Accordingly the
ledger closes their normalization steps while retaining those exact
existence suppliers.

\medskip
\noindent\textbf{Source and scope record.}
The companion provenance ledger gives declaration-level pinned Git
ranges and hashes for every SE claim. SE1--SE6 follow the substantive
\texttt{LatticeEmbeddingMemUnitGroup} proof; SE7--SE8 reconstruct
\texttt{fits} and \texttt{ModelLamPDef}; SE9--SE10 supply
\texttt{ExistsInnerProductSpace}; SE11 proves the four-field
\texttt{CDHaar} conclusion by an ordinary Lie-coordinate argument;
SE12 supplies the exact bare \texttt{GJStage}; SE13--SE14 leave the
full \texttt{NStage}/\texttt{LocStage} supplier open; SE15--SE16
reconstruct the coordinate and scalar pin normalization. None of
these constructions identifies an auxiliary stage with an externally
specified geometric or spectral realization without the relevant tie.
